% SPDX-License-Identifier: CC-BY-SA-4.0
% Author: Matthieu Perrin
% Part: <Nom de la partie>
% Section: <Nom de la section>
% Sub-section: <Nom de la sous-section>  % (facultatif, laisser vide si non utilisé)
% Frame: <Titre de la slide>

\begingroup

\begin{frame}{Indécidabilité de l'ambiguïté}

  \onBlock[top=-5mm]{Théorème -- Indécidabilité de l'ambiguïté}{
    Le problème \textsc{Ambiguity} est indécidable
  }
  
  \onBlock[y=7mm]{Démonstration}{
    La fonction $\varphi$ suivante est une réduction de \textsc{pcp} vers \textsc{Ambiguity} :
    {\small
      $$
      \varphi : \left\{\begin{array}{@{}r@{~}c@{~}l@{}}
      \mathcal{I}(\textsc{pcp})
      &
      \rightarrow
      &
      \mathcal{I}(\textsc{Ambiguity}) \text{\quad\quad\quad\example{(sur l'alphabet $\Sigma \cup D$})}
      %
      \vspace{2mm}\\
      %
      \left\{\VDomino{\alpha_1}{\beta_1}, \ldots, \VDomino{\alpha_n}{\beta_n} \right\}
      &
      \mapsto
      &
      \left\{\begin{array}{@{~}r@{~~}c@{~~}l@{}l@{~}c@{~}l@{~}l@{~}c@{~}l@{~}l@{~}c@{~}l@{~}l@{}}
      S &\rightarrow& A \mid B
      \vspace{1mm}\\
      %
      A &\rightarrow&
      \alpha_1 A & \VDomino{\alpha_1}{\beta_1}
      &\mid&
      \alpha_1   & \VDomino{\alpha_1}{\beta_1}
      & \mid \ldots \mid &
      \alpha_n A & \VDomino{\alpha_n}{\beta_n}
      & \mid &
      \alpha_n   & \VDomino{\alpha_n}{\beta_n}
      \vspace{1mm}\\
      %
      B &\rightarrow&
      \beta_1 B  & \VDomino{\alpha_1}{\beta_1}
      &\mid&
      \beta_1    & \VDomino{\alpha_1}{\beta_1}
      & \mid \ldots \mid &
      \beta_n B  & \VDomino{\alpha_n}{\beta_n}
      & \mid &
      \beta_n    & \VDomino{\alpha_n}{\beta_n} 
      \end{array}\right.
      \end{array}\right.
      $$
    }
  }

  \onExampleBlock[bottom=-3mm]{Exemple}{
    \begin{itemize}
    \item\vspace{-2mm} $D = {\color{black}\left\{{\footnotesize \VDomino[.7]{a}{baa}, \VDomino[.7]{ab}{aa}, \VDomino[.7]{bba}{bb}} \right\}}$
      a pour solution ${\footnotesize \VDomino{bba}{bb}\, \VDomino{ab}{aa}\, \VDomino{bba}{bb}\, \VDomino{a}{baa}}$.
    \item $\varphi(D)$ a deux dérivations pour : 
      $bbaabbbaa\, \scriptsize
      \VDomino{a}{baa}\, 
      \VDomino{bba}{bb}\,
      \VDomino{ab}{aa}\,
      \VDomino{bba}{bb}
      $.
      \begin{itemize}
      \item $S \vdash A \vdash^\star \structure{bba}\cdot \structure{ab} \cdot \structure{bba} \cdot \structure{a} \cdot \scriptsize
        \VDomino{a}{baa}\cdot
        \VDomino{bba}{bb}\cdot
        \VDomino{ab}{aa}\cdot
        \VDomino{bba}{bb}$
      \item $S \vdash B \vdash^\star \example{bb}\cdot \example{aa} \cdot \example{bb} \cdot \example{baa} \cdot\hspace{-.2mm} \scriptsize
        \VDomino{a}{baa}\cdot
        \VDomino{bba}{bb}\cdot
        \VDomino{ab}{aa}\cdot
        \VDomino{bba}{bb}$
      \end{itemize}
    \end{itemize}
  }

\end{frame}

\endgroup
\endinput
